\section{Predictive States and Planning Geometry}
\label{sec:framework}
\label{sec:method}

\paragraph{Language-mediated control.}
Language provides many surface forms for the same situation. ``The number of red objects is six'' and ``There are six red objects'' may support exactly the same future calculations. Replacing either sentence by ``There are not six red objects'' can change every valid continuation. A useful state representation should ignore the first wording change and preserve the second semantic change.

We model this setting as a controlled process
\begin{equation}
    \mathcal{M}=(\mathcal{S},\mathcal{U},P,\Omega,O,c,\mathcal{G}).
\end{equation}
The semantic state is $s_t\in\mathcal{S}$ and a semantic intervention is $u_t\in\mathcal{U}$. The agent receives a language observation $x_t\in\Omega$ and chooses a language action $a_t$. A grounding map $\alpha(a_t)=u_t$ connects the phrase to its effect. Decisions may depend on the full history
\begin{equation}
    h_t=(x_0,a_0,x_1,\ldots,a_{t-1},x_t),
\end{equation}
rather than on the latest sentence alone. Histories accommodate synonymous descriptions, irrelevant facts, and partial observations without requiring a network to recover a privileged hidden coordinate system.

\paragraph{Predictive equivalence.}
Predictive state representations define state through action-conditioned predictions of future observations \citep{littman2001predictive}. Let $Y$ collect future task-relevant observations and terminal outcomes, let $C$ collect future costs, and let $F$ collect future feasible-action sets. Histories $h$ and $h'$ are controlled-predictively equivalent when every policy over actions admissible from both histories induces the same future law:
\begin{align}
h \sim_{\mathrm{pred}} h'
\quad\Longleftrightarrow\quad
\mathcal{L}_{\pi}(Y_{t+1:\tau},C_{t:\tau},F_{t:\tau}\mid h)
=
\mathcal{L}_{\pi}(Y_{t+1:\tau},C_{t:\tau},F_{t:\tau}\mid h')
\quad\forall\pi.
\label{eq:predictive-equivalence}
\end{align}
The corresponding predictive state is the equivalence class
\begin{equation}
    q(h)=[h]_{\sim_{\mathrm{pred}}}\in\mathcal{H}/{\sim_{\mathrm{pred}}}.
\end{equation}
Histories with identical controlled futures meet in the quotient even when their wording differs. A negation, operator substitution, or changed entity relation remains visible whenever it changes a future consequence. Information that cannot affect any future task outcome may disappear.

\begin{proposition}[Control sufficiency]
If $h\sim_{\mathrm{pred}}h'$ includes equality of future costs and terminal events under every admissible policy, then $V^\pi(h)=V^\pi(h')$ for every policy $\pi$, and therefore $V^\star(h)=V^\star(h')$.
\end{proposition}
\begin{proof}
Return is a function of the future cost and terminal sequence. Equation~\ref{eq:predictive-equivalence} assigns that sequence the same distribution under each policy. Its expectation is therefore equal. Maximizing over policies preserves the equality.
\end{proof}

Proposition 1 defines the state required for control. We do not claim it as a new control theorem. Bisimulation gives a related recursive characterization through agreement in immediate reward and transitions between equivalence classes \citep{ferns2011bisimulation}. GAR imposes a narrower condition, namely goal-conditioned order among actions available at one history.

The experiments instantiate the framework with a deterministic finite-horizon process. Executing $a$ appends one deterministic textual consequence, written $T(h,a)$. Stochastic environments would require a predictive distribution or another representation of multimodal futures. A point predictor trained by mean-squared error generally represents only a conditional mean.

\paragraph{Coordinate ambiguity.}
Suppose an encoder $\phi$ and predictor $F$ recover every deterministic transition:
\begin{equation}
    F(\phi(h),e(a))=\phi(T(h,a)).
\end{equation}

\begin{proposition}[Predictive ambiguity]
Let $b$ be a bijection on the realized latent states and define $\widetilde\phi=b\circ\phi$ together with
\begin{equation}
    \widetilde F(b(z),e(a))=b(F(z,e(a))).
\end{equation}
Then $(\widetilde\phi,\widetilde F)$ predicts the same realized transitions exactly. A non-isometric $b$ need not preserve the order of successor distances to a goal.
\end{proposition}
\begin{proof}
Substitution gives $\widetilde F(\widetilde\phi(h),e(a))=\widetilde\phi(T(h,a))$. On a finite realized state set, the images of two successors and a goal may be placed at arbitrary distinct points. Their goal-distance order can change without changing any transition identity.
\end{proof}

The two embeddings in Figure~\ref{fig:nonidentifiability} realize the same transitions and disagree on the preferred successor. Additional assumptions can reduce this ambiguity. For example, controlled Gaussian world models with spectral separation and sufficient conditional action excitation can be identified up to an orthogonal map \citep{song2026identifiability}. Proposition 2 concerns predictive learning without such metric-preserving guarantees.

\paragraph{Predictive architecture.}
The online encoder produces $z_h=\phi_\theta(h)$ and the action encoder produces $e_\theta(a)$. Their predicted successor is
\begin{equation}
    \widehat z_{h,a}=F_\theta(z_h,e_\theta(a)).
\end{equation}
An exponential-moving-average encoder maps each observed successor history to a stopped-gradient target $\bar z_{T(h,a)}$. Factual and counterfactual prediction losses match $\widehat z_{h,a}$ to these targets. The target encoder stays in evaluation mode and every dropout rate is zero.

At inference time, the prompt and query are represented by a problem anchor $z_x$. The transition-value realization assigns a lower-is-better energy
\begin{equation}
    E^{\mathrm{TV}}_\psi(h,a,x)=V_\psi(\widehat z_{h,a},z_x).
\end{equation}
The policy predicts and scores every currently feasible phrase, then selects the minimum-energy candidate. Figure~\ref{fig:method} separates this inference path from the training signals.

\begin{figure}[t]
    \centering
    \resizebox{0.96\linewidth}{!}{\begin{tikzpicture}[
    x=1cm,y=1cm,
    box/.style={draw=black!65,rounded corners=2pt,fill=white,minimum height=7mm,align=center,font=\scriptsize,inner xsep=5pt},
    latent/.style={circle,draw=blue!60!black,fill=blue!7,minimum size=7mm,font=\scriptsize},
    target/.style={circle,draw=green!50!black,fill=green!8,minimum size=7mm,font=\scriptsize},
    arrow/.style={-{Latex[length=2mm]},thick},
    targetarrow/.style={-{Latex[length=2mm]},thick,green!45!black,dashed},
    control/.style={-{Latex[length=2mm]},thick,black!55,densely dashed}
]
% Deployed predictive policy path
\node[box] (history) at (0,2.15) {history $h_t$\\problem anchor $x$};
\node[box] (encoder) at (1.75,2.15) {online encoder\\$\phi_\theta$};
\node[latent] (state) at (3.05,2.15) {$z_t$};
\node[box] (actions) at (0,0.95) {feasible phrases\\$a_1,\ldots,a_K$};
\node[box] (aenc) at (1.75,0.95) {action encoder\\$e_\theta$};
\node[box] (pred) at (4.45,1.55) {predict successor\\$\hat z_i=F_\theta(z_t,e(a_i))$};
\node[box] (value) at (6.65,1.55) {score consequence\\$E_i=V_\psi(\hat z_i,z_x)$};
\node[box] (choice) at (8.55,1.55) {select phrase\\$\arg\min_i E_i$};
\draw[arrow] (history) -- (encoder);
\draw[arrow] (encoder) -- (state);
\draw[arrow] (state) -- (pred);
\draw[arrow] (actions) -- (aenc);
\draw[arrow] (aenc) -- (pred);
\draw[arrow] (pred) -- (value);
\draw[arrow] (value) -- (choice);

% Training target path
\node[box] (outcomes) at (0,-0.55) {observed alternative\\outcomes};
\node[box] (ema) at (1.75,-0.55) {EMA encoder\\$\bar\phi$};
\node[target] (next) at (3.05,-0.55) {$\bar z_i$};
\node[target] (goal) at (4.05,-0.55) {$\bar z_g$};
\node[box,draw=red!60!black,fill=red!3] (gar) at (5.6,-0.55) {GAR target\\order + gap};
\node[box,draw=red!60!black,fill=red!3] (loss) at (7.35,-0.55) {$\mathcal L_{\rm rank}$\\$+\mathcal L_{\rm reg}$};
\draw[arrow] (outcomes) -- (ema);
\draw[targetarrow] (ema) -- node[above,font=\tiny] {stop grad.} (next);
\draw[targetarrow] (next) -- (gar);
\draw[targetarrow] (goal) -- (gar);
\draw[arrow,red!65!black] (gar) -- (loss);
\draw[arrow,red!65!black] (value) -- (loss);

% Information-matched architectural control
\node[box,draw=black!45,fill=black!2] (direct) at (9.15,-0.55) {direct GAR\\$D_\psi(z_t,e(a_i),z_x)$};
\draw[control] (gar) -- (direct);
\draw[control] (direct) -- (choice);
\end{tikzpicture}
}
    \caption{Inference and training paths for predictive GAR. At inference, each candidate phrase is encoded, passed through the successor predictor, and scored against the problem anchor. Training observes factual and counterfactual successor histories. EMA target distances define within-state order labels, while latent prediction preserves the action-conditioned transition.}
    \label{fig:method}
\end{figure}

\paragraph{Geometric Advantage Ranking.}
Let $\bar z_g$ be the EMA representation of the terminal training history. Each observed candidate successor receives the target distance
\begin{equation}
    \bar d_i=d(\bar z_{T(h,a_i)},\bar z_g).
\end{equation}
The default target evaluates a candidate after a bounded two-step reference continuation. It remains a geometric progress surrogate rather than a Bellman-consistent advantage. Whether its order agrees with true continuation quality is measured directly in the action audits.

Whenever $\bar d_i+\epsilon<\bar d_j$, GAR imposes
\begin{equation}
    \mathcal{L}_{\mathrm{rank}}
    =\sum_{i\succ j}\left[m+E_\psi(h,a_i,x)-E_\psi(h,a_j,x)\right]_+.
\end{equation}
A calibration term matches differences in predicted energy to differences in target distance:
\begin{equation}
    \mathcal{L}_{\mathrm{reg}}
    =\sum_{i<j}\left[(E_j-E_i)-\beta(\bar d_j-\bar d_i)\right]^2.
\end{equation}
Together with factual and counterfactual latent prediction, the training objective is
\begin{equation}
    \mathcal{L}=\mathcal{L}_{\mathrm{factual}}
    +\lambda_{\mathrm{cf}}\mathcal{L}_{\mathrm{cf}}
    +\lambda_{\mathrm{rank}}\mathcal{L}_{\mathrm{rank}}
    +\lambda_{\mathrm{reg}}\mathcal{L}_{\mathrm{reg}}.
\end{equation}
GAR gradients reach the scorer, predictor, action encoder, and online history encoder in the default model. The EMA network follows the online encoder without receiving gradients.

GAR names the supervision, not the scorer architecture. Transition-value GAR factors the energy through $\widehat z_{h,a}$. Direct GAR uses $D_\psi(z_h,e(a),z_x)$ and bypasses the successor in the policy path. Raw-distance GAR uses $d(\widehat z_{h,a},\bar z_g)$ without a learned head. The last variant receives the terminal training representation during evaluation and is reported only as an oracle-terminal diagnostic.

\paragraph{Local planning condition.}
Let $D^\star(h)$ denote the shortest remaining action count in a deterministic unit-cost puzzle, and define
\begin{equation}
    A^\star(h,a)=D^\star(h)-1-D^\star(T(h,a)).
\end{equation}
An action on a shortest solution has $A^\star=0$. An action that spends a step without reducing the remaining distance has $A^\star=-1$.

\begin{proposition}[Local ordering]
If a learned score ranks every optimal action above every suboptimal action at every reachable history, greedy selection reaches the goal in exactly $D^\star(h_0)$ actions.
\end{proposition}
\begin{proof}
At remaining distance $d>0$, the selected action reaches a history with distance $d-1$. Repeating the argument reaches distance zero after $d$ actions.
\end{proof}

If the smallest true action gap is $\gamma>0$, a uniform estimation error below $\gamma/2$ preserves the same ordering. GAR does not receive $A^\star$ during main training. The proposition instead identifies the behavioral condition that its geometric labels must recover. Pairwise ordering, top-1 recall, regret, and margin therefore provide more direct tests than effective rank or global probe accuracy.

\paragraph{Compression.}
The quotient is the coarsest exact partition that preserves controlled futures. It supplies a formal notion of task-relevant compression, but the neural encoder is not guaranteed to find that minimum. The information bottleneck likewise seeks a compressed code that preserves target information \citep{tishby2000information}. Our objective contains no rate penalty, so latent width and effective rank cannot be interpreted as evidence for an explicit information bottleneck.
